% =============================================================================
% ЛЕКЦИЯ 1
% =============================================================================

\lecture{1}{Введение в математический анализ}

\lecturedate{1 сентября 2025}

\section{Вещественные числа}

Множество вещественных чисел $\R$ является фундаментальным объектом математического анализа. Оно обладает структурой \term{полного упорядоченного поля}.

\subsection{Аксиомы поля}

Для любых $a, b, c \in \R$ выполняются следующие свойства:

\begin{enum}
    \item Коммутативность: $a + b = b + a$, \quad $ab = ba$
    \item Ассоциативность: $(a+b)+c = a+(b+c)$
    \item Дистрибутивность: $a(b+c) = ab + ac$
    \item Существование нейтральных элементов: $0$ и $1$
    \item Существование обратных элементов
\end{enum}

\subsection{Аксиома полноты}

\begin{defn}[Супремум]
Пусть $A \subset \R$ — непустое ограниченное сверху множество. \term{Супремумом} множества $A$ называется наименьшая верхняя грань:
\[
    \sup A = \min\set{M \in \R : \forall a \in A,\, a \leq M}
\]
\end{defn}

\begin{thm}[О супремуме]
Всякое непустое ограниченное сверху множество $A \subset \R$ имеет супремум.
\end{thm}

\pf{
Это свойство принимается как аксиома полноты вещественных чисел и не может быть доказано из остальных аксиом.
}

\section{Последовательности}

\begin{defn}[Последовательность]
\term{Последовательность} — это функция $a: \N \to \R$, обозначаемая $(a_n)_{n=1}^{\infty}$ или просто $(a_n)$.
\end{defn}

\subsection{Сходимость последовательностей}

\begin{defn}[Предел последовательности]
Число $L \in \R$ называется \term{пределом} последовательности $(a_n)$, если:
\[
    \forall \eps > 0 \;\; \exists N \in \N: \;\; \forall n > N \then \abs{a_n - L} < \eps
\]
Обозначение: $\limn a_n = L$ или $a_n \to L$.
\end{defn}

\begin{thm}[Единственность предела]
Если последовательность $(a_n)$ сходится, то её предел единственен.
\end{thm}

\pf{
Предположим, что $a_n \to L_1$ и $a_n \to L_2$, где $L_1 \neq L_2$. Положим $\eps = \frac{\abs{L_1 - L_2}}{2} > 0$.

По определению предела:
\begin{items}
    \item $\exists N_1: \forall n > N_1 \then \abs{a_n - L_1} < \eps$
    \item $\exists N_2: \forall n > N_2 \then \abs{a_n - L_2} < \eps$
\end{items}

При $n > \max(N_1, N_2)$:
\[
    \abs{L_1 - L_2} \leq \abs{a_n - L_1} + \abs{a_n - L_2} < 2\eps = \abs{L_1 - L_2}
\]
Противоречие. Следовательно, $L_1 = L_2$.
}

\section{Примеры}

\begin{example}
Докажем, что $\limn \frac{1}{n} = 0$.

Пусть $\eps > 0$. Выберем $N = \ceil{\frac{1}{\eps}}$. Тогда для всех $n > N$:
\[
    \abs{\frac{1}{n} - 0} = \frac{1}{n} < \frac{1}{N} \leq \eps
\]
\end{example}

\important{Понимание определения предела через $\eps$-$N$ — ключевой навык в математическом анализе.}

\begin{remark}
На практике часто используют арифметику пределов вместо прямой проверки определения.
\end{remark}
